\chapter{Command-line shortcuts}\label{ch:cmdline}

You can give ProPrep extra instructions when you start it, by typing
\emph{options} after the command name in the Terminal. Options let you skip
start-up questions, load a structure straight away, replay a recorded session
or run a batch of proteins. This chapter lists the options a user is likely to
need, grouped by purpose, with an example of each. None of them is required:
everything can also be done from the menus.

\section{How to type options}\label{sec:cmdline-howto}

\begin{itemize}
  \item Each option starts with two hyphens, such as \typed{--full-menu}, and is
  separated from the next by a space. The order does not matter.
  \item Some options need a value after them: a PDB code, a file name or a
  folder, for example \typed{--pdbid 1UBQ}.
  \item Put text that contains spaces in quotation marks:
  \typed{--session-description "first try"}.
  \item For folders, give the full path. Starting it with \typed{\textasciitilde{}/} (your home
  folder) is the easiest way; the pitfalls in Section~\ref{sec:cmdline-pitfalls}
  explain why.
  \item \typed{proprep --help} prints the complete list and exits. So do
  \typed{proprep-web --help} and the library command in
  Section~\ref{sec:cmdline-library}.
\end{itemize}

You can combine several options in one command:

\begin{terminal}
proprep --full-menu --pdbid 1CRN --output-dir ~/ProPrep-work/crambin
\end{terminal}

\section{Starting and menus}\label{sec:cmdline-start}

\begin{longtable}{@{}p{0.33\linewidth}p{0.62\linewidth}@{}}
\toprule
\textbf{Option} & \textbf{What it does, with an example} \\
\midrule
\endhead
\typed{--full-menu} & Start in the Full Menu, the menu used in this manual;
skips the menu question.\newline
\typed{proprep --full-menu} \\
\typed{--workflow} & Start in Workflow Mode; skips the menu question.\newline
\typed{proprep --workflow} \\
\typed{--set-default} \textit{mode} & Save the menu ProPrep starts in, then
exit. \textit{mode} is \typed{full-menu}, \typed{workflow} or \typed{ask}.\newline
\typed{proprep --set-default full-menu} \\
\typed{--menu-grid} & Draw the Full Menu as two columns of panels, for this run
only.\newline
\typed{proprep --full-menu --menu-grid} \\
\typed{--menu-list} & Draw the Full Menu as one column (the normal layout), for
this run only. \\
\typed{--version} & Print the version and exit.\newline
\typed{proprep --version} gives
\texttt{ProPrep v1.23.0 - Proper Protein Preparation workflow manager}. \\
\typed{--help}, \typed{-h} & List every option and exit. \\
\bottomrule
\end{longtable}

\section{Project folder and structure}\label{sec:cmdline-project}

\begin{longtable}{@{}p{0.33\linewidth}p{0.62\linewidth}@{}}
\toprule
\textbf{Option} & \textbf{What it does, with an example} \\
\midrule
\endhead
\typed{--output-dir} \textit{folder} & Use this folder as the project folder,
creating it (and any folders above it) if needed. Skips the project-name
question. Give a full path.\newline
\typed{proprep --full-menu --output-dir \textasciitilde{}/ProPrep-work/ubq} \\
\typed{--pdbid} \textit{code} & Download this structure from the Protein Data
Bank at start-up. Skips the project-name question and uses the current
folder.\newline
\typed{proprep --full-menu --pdbid 1UBQ} \\
\typed{--pdbfile} \textit{file} & Load a PDB file from your computer at
start-up. Skips the project-name question.\newline
\typed{proprep --full-menu --pdbfile my\_protein.pdb} \\
\bottomrule
\end{longtable}

If the download fails, ProPrep says so and opens the main menu with nothing
loaded. This happened with \typed{--pdbid ZZZZ}, a code that does not exist. If you give both \typed{--pdbid} and
\typed{--pdbfile}, ProPrep uses \typed{--pdbid} and ignores the file without a
message.

\section{Sessions}\label{sec:cmdline-sessions}

Chapter~\ref{ch:sessions} explains sessions, replay and the session editor.

\begin{longtable}{@{}p{0.33\linewidth}p{0.62\linewidth}@{}}
\toprule
\textbf{Option} & \textbf{What it does, with an example} \\
\midrule
\endhead
\typed{--no-session} & Do not record this run. No session file is written,
ProPrep does not offer to resume, and \key{undo} is not available. Same as
\typed{--no-auto-record}.\newline
\typed{proprep --full-menu --no-session} \\
\typed{--session-description} \textit{text} & Store a note in the session
file; it appears in the Notes column when you choose a session to resume.\newline
\typed{proprep --full-menu --session-description "1UBQ, pH 7"} \\
\typed{--session-dir} \textit{folder} & Where to write the session file and look
for old ones. The folder must already exist; give a full path.\newline
\typed{proprep --full-menu --session-dir \textasciitilde{}/ProPrep-work/sessions} \\
\typed{--resume-session} \textit{file} & Open a recorded session to continue
it, edit it, or cancel and start fresh.\newline
\typed{proprep --full-menu --resume-session proprep\_session\_20261121\_130253.json} \\
\typed{--force-replay} \textit{file} & Play a recorded session back exactly,
without your input, and stop at its end.\newline
\typed{proprep --full-menu --force-replay proprep\_session\_20261121\_122825.json} \\
\typed{--demo-delay} \textit{seconds} & With \typed{--force-replay} or
\typed{--resume-session}: wait this long between replayed answers, so you can
watch.\newline
\typed{... --force-replay session.json --demo-delay 0.5} \\
\bottomrule
\end{longtable}

\section{Templates and batch processing}\label{sec:cmdline-batch}

Chapter~\ref{ch:batch} explains these step by step.

\begin{longtable}{@{}p{0.33\linewidth}p{0.62\linewidth}@{}}
\toprule
\textbf{Option} & \textbf{What it does, with an example} \\
\midrule
\endhead
\typed{--create-template} \textit{session} & Turn a recorded session into a
template, marking the answers (such as the PDB code) that change from protein
to protein.\newline
\typed{proprep --create-template ubq\_session.json} \\
\typed{--template-info} \textit{template} & Show a template's variables and
size.\newline
\typed{proprep --template-info ubq\_session\_template.json} \\
\typed{--validate-template} \textit{template} & Check a template without running
it. Add \typed{--batch-list} to check a list of proteins too.\newline
\typed{proprep --validate-template ubq\_session\_template.json --batch-list ubq\_session\_batch.csv} \\
\typed{--batch-replay} \textit{template} & Run the template on every protein in
the list given with \typed{--batch-list}. \\
\typed{--batch-list} \textit{file} & The list of proteins: a CSV file with
variable names in the first row, or a plain list with one PDB code per line. \\
\typed{--batch-continue} & Keep going when one protein fails, instead of
stopping. \\
\typed{--output-dir} \textit{folder} & With \typed{--batch-replay}: where the
run folders \file{run_01}, \file{run_02}, \ldots{} and the report go. The folder
must already exist. \\
\typed{--web}, \typed{--no-web} & Replay in browser mode or in plain Terminal
mode. Normally ProPrep picks the mode the template was recorded in. \\
\bottomrule
\end{longtable}

A complete batch command looks like this:

\begin{terminal}
mkdir batch_out
proprep --batch-replay ubq_session_template.json \
        --batch-list ubq_session_batch.csv \
        --output-dir batch_out --batch-continue
\end{terminal}

The \texttt{\textbackslash} at the end of a line lets you continue a long
command on the next line; you can also type it all on one line without them.

\section{Viewer and analysis shortcuts}\label{sec:cmdline-shortcuts}

\begin{longtable}{@{}p{0.33\linewidth}p{0.62\linewidth}@{}}
\toprule
\textbf{Option} & \textbf{What it does, with an example} \\
\midrule
\endhead
\typed{--pdbview} \textit{code or file} & Load a structure and open it in the
3D viewer straight away (at an address such as
\url{http://localhost:8766/viewer}), then show the main menu. Nothing is
recorded. Add \typed{--full-menu} to skip the menu question.\newline
\typed{proprep --full-menu --pdbview 1UBQ} \\
\typed{--analysis} & Open the simulation analysis file browser straight away
(Chapter~\ref{ch:analysis}). Start it from inside the project folder. Cannot be
combined with \typed{--pdbid} or \typed{--pdbfile}.\newline
\typed{proprep --full-menu --analysis} \\
\bottomrule
\end{longtable}

\section{When something goes wrong}\label{sec:cmdline-debug}

\begin{longtable}{@{}p{0.33\linewidth}p{0.62\linewidth}@{}}
\toprule
\textbf{Option} & \textbf{What it does, with an example} \\
\midrule
\endhead
\typed{--debug} & Show extra messages, and the full technical error report if
ProPrep stops with an error.\newline
\typed{proprep --full-menu --debug} \\
\typed{--verbose}, \typed{-v} & Fill the screen with detailed technical log
lines. Useful only when the developers ask for it.\newline
\typed{proprep --full-menu --verbose} \\
\bottomrule
\end{longtable}

Chapter~\ref{ch:help} says more about both. Two further options are tools for
the ProPrep developers, and users do not need them:
\typed{--generate-feedback-keypair} and \typed{--decrypt-feedback}.

\section{Options for the browser version}\label{sec:cmdline-web}

These go after \typed{proprep-web}.

\begin{longtable}{@{}p{0.33\linewidth}p{0.62\linewidth}@{}}
\toprule
\textbf{Option} & \textbf{What it does, with an example} \\
\midrule
\endhead
\textit{theme} (a word) & Color scheme of the terminal pane: \typed{dark}
(default), \typed{light}, \typed{black} or \typed{white}.\newline
\typed{proprep-web white} \\
\typed{--theme} \textit{theme} & The same, written as an option. \\
\typed{--font-size} \textit{pixels} & Text size in the terminal pane, 8 to 32
(default 13). You can also change it in the page.\newline
\typed{proprep-web --font-size 16} \\
\typed{--high-contrast} & Stronger text contrast, for easier reading. \\
\typed{--port} \textit{number} & The port to start from (default 8000). If it is
busy, the next free one is used.\newline
\typed{proprep-web --port 8420} \\
\typed{--strict-port} & Stop with an error instead of trying the next port. \\
\typed{--no-browser} & Do not open the browser; open the printed address
yourself. \\
\typed{--no-transcript} & Do not keep a plain-text log of the session in
\file{~/.proprep/web_sessions/}. \\
\typed{--host}, \typed{--log-level} & Technical settings for the web server;
leave them alone unless you know you need them. \\
\typed{--seats} \textit{N}, \typed{--seats-dir}, \typed{--public-url} &
Workshop mode: one ProPrep for several attendees, each with a private link.
Tried only for starting and stopping in this manual. \\
\bottomrule
\end{longtable}

\section{The parameter library command}\label{sec:cmdline-library}

ProPrep keeps a personal library of force-field parameters that you made
(Chapter~\ref{ch:forcefields}). A separate command, \typed{proprep-library},
lets you look at it, save copies (\term{snapshots}) of it, empty it and bring
an old copy back. It never deletes anything: emptying or restoring first saves
a snapshot of the current library.

\begin{longtable}{@{}p{0.33\linewidth}p{0.62\linewidth}@{}}
\toprule
\textbf{Command} & \textbf{What it does} \\
\midrule
\endhead
\typed{status} & How many files the library holds, and how many snapshots
exist. \\
\typed{list} & List the snapshots with date, contents and note. \\
\typed{snapshot} \textit{name} \typed{-m} \textit{note} & Save a copy of the
library now. Without a name, the date and time are used. \\
\typed{reset} & Save a snapshot, then empty the library. Asks first; \typed{-y}
skips the question. \\
\typed{restore} \textit{name} & Save a snapshot of the current library, then
replace it with the named snapshot. \\
\typed{--root} \textit{folder} & Work on a different library folder instead of
\file{~/.proprep}. \\
\bottomrule
\end{longtable}

\begin{watchout}
In the copy installed for this manual, the \typed{proprep-library} command did
not start; it stopped with \texttt{bad interpreter: ... python3.13: no such file
or directory}. The same tool works when started through Python. With the
ProPrep environment active (Chapter~\ref{ch:installation}), type
\typed{python -m proprep.utils.user\_library\_snapshots} followed by the
command, for example:
\begin{terminal}
python -m proprep.utils.user_library_snapshots status
python -m proprep.utils.user_library_snapshots snapshot before_test -m "my note"
\end{terminal}
\end{watchout}

\section{Pitfalls}\label{sec:cmdline-pitfalls}

\begin{watchout}
\textbf{\typed{--output-dir} with a short (relative) path.} With
\typed{--output-dir a/b}, ProPrep puts the files in the right place but stores a
doubled path, \file{.../a/b/a/b}, as the project folder. Tools that use that
stored path may then look in the wrong place. Give a full path instead, for
example \typed{--output-dir \textasciitilde{}/ProPrep-work/a/b}.
\end{watchout}

\begin{watchout}
\textbf{\typed{--session-dir} must be a full path to a folder that exists.}
With a relative path, or a folder that does not exist yet, ProPrep stops right
after the project question with \texttt{Error: [Errno 2] No such file or
directory}. Create the folder first (\typed{mkdir}) and give its full path.
\end{watchout}

\begin{watchout}
\textbf{The batch output folder must exist.} With \typed{--batch-replay}, an
\typed{--output-dir} that does not exist makes every run fail with
\texttt{Setup error: [Errno 2] No such file or directory}, followed by a long
technical error report. Create it first with \typed{mkdir}.
\end{watchout}

\begin{watchout}
\textbf{\typed{--pdbview} still asks which menu you want} unless you also give
\typed{--full-menu} (or \typed{--workflow}). The same is true of
\typed{--force-replay}, because the menu choice is not part of the recording.
\end{watchout}

\begin{itemize}\raggedright
  \item Options that work only together with another one say so and are
  ignored, for example
  \texttt{Warning: --demo-delay has no effect without --force-replay or
  --resume-session}.
  \item Some options cannot be combined. \typed{--workflow} with
  \typed{--full-menu}, or \typed{--menu-grid} with \typed{--menu-list}, gives an
  error such as
  \texttt{proprep: error: argument --full-menu: not allowed with argument
  --workflow}.
  \item \typed{proprep-web --strict-port} on a busy port first prints its
  \texttt{listening on} line and only then the error
  \texttt{address already in use}. Trust the error.
\end{itemize}
